This activity covers user stories US-025 and US-027. The analysis component derives two kinds of results from the stored data: descriptive metrics of the dataset and a sentiment label for each comment.

\subsection{Descriptive metrics}

The descriptive metrics of a project are the number of collection runs, distinct videos, distinct channels, top-level comments, replies and statistics snapshots, together with the period in which the project's runs discovered its videos. They are not stored in a table: as anticipated in Activity~\ref{act:1.3}, they are computed on request by PostgreSQL through Django aggregates, using distinct counts and filtered counts, so they always reflect the data that remain after the purge and no row is loaded into Python to count it.

Because videos, channels and comments are shared between projects (Activity~\ref{act:3.2}), every metric is scoped through the project's own collection runs. Videos and channels are counted through the run links, so a video found by two runs counts once, and comments, replies and snapshots are counted through the project's runs, never through the videos, so another researcher's comments on a shared video are not included. The collection period is taken from the moment the project's runs discovered each video, not from the video's last retrieval time, which any other run can overwrite and which would reveal another researcher's activity. The counters stored on each collection run and processing run, described in Activities~\ref{act:2.1} and~\ref{act:3.1}, complement these metrics because they survive the purge.

\subsection{Sentiment classification}

Sentiment analysis is implemented against a contract rather than against a specific model. A classifier is any component that exposes a model name, a model version and a \texttt{classify} method that receives a list of texts, since models are faster in batches, and returns one result per text, with a label, a score between $-1$ and 1, a confidence between 0 and 1, a language code and the raw model output. The model is selected with a configuration value, \texttt{SENTIMENT\_\allowbreak MODEL}; an empty value disables the step, and a value that does not name a registered classifier stops the application at startup with a Django system check error, so a typing error cannot silently skip the analysis.

The classification runs as its own processing run, after the normalization run of the same collection completes (Activity~\ref{act:3.1}). It is not run after a failed normalization, because shared comments could still hold text normalized by an earlier run under older rules, and classifying that mixture would be misleading. Comments without normalized text are counted as skipped. Every result is validated before it is stored: the label must be one of \texttt{positive}, \texttt{negative}, \texttt{neutral}, \texttt{mixed} or \texttt{unknown}, the score and confidence must be numbers in their ranges, Boolean and not-a-number values are rejected and the classifier must return exactly one result per text. An invalid result fails the processing run without storing it, and PostgreSQL check constraints enforce the same ranges and labels as a second barrier.

Each stored result carries the model name and version, and the report of the processing run records the model, the number of inputs, the count per label and the number of skipped comments, without text or identifiers. Results of different models or model versions are therefore never mixed, and changing the model produces a new, separately reported pass instead of overwriting the previous one.

% TODO(US-025 part B): name the pre-trained sentiment model registered in CLASSIFIERS, its version, licence and label mapping, and the reason for choosing it. Until then no real model is registered, and only the test classifier produces results.
